\documentclass{article}
\usepackage[utf8]{inputenc}
\usepackage[T1]{fontenc}
\usepackage{hyperref}
\usepackage{listings}
\usepackage{geometry}
\geometry{a4paper, margin=1in}

\title{Medical Imaging Project: Kidney Analysis with Multi-Phase CT Scans}
\author{Project Documentation}
\date{\today}

\begin{document}

\maketitle

\section{Project Overview}

This project focuses on advanced kidney analysis using multi-phase CT (Computed Tomography) imaging data. The study involves comprehensive segmentation, measurement, and functional assessment of kidney function through contrast-enhanced CT scans and eGFR (estimated Glomerular Filtration Rate) calculations.

\section{CT Scan Data with Different Contrast Phases}

\subsection{Multi-Phase CT Imaging Protocol}

Multi-phase CT imaging is a sophisticated technique that captures images at different time points after contrast agent injection, allowing for comprehensive assessment of kidney function and anatomy.

\subsubsection{Arterial Phase (Early Phase)}
\begin{itemize}
    \item \textbf{Timing}: 15-25 seconds after contrast injection
    \item \textbf{Purpose}: Captures the arterial blood supply to the kidneys
    \item \textbf{Key Features}:
    \begin{itemize}
        \item High contrast in renal arteries
        \item Clear visualization of arterial anatomy
        \item Assessment of arterial stenosis or abnormalities
        \item Early enhancement of kidney cortex
    \end{itemize}
\end{itemize}

\subsubsection{Venous Phase (Portal Phase)}
\begin{itemize}
    \item \textbf{Timing}: 60-90 seconds after contrast injection
    \item \textbf{Purpose}: Captures venous drainage and parenchymal enhancement
    \item \textbf{Key Features}:
    \begin{itemize}
        \item Peak enhancement of kidney parenchyma
        \item Visualization of renal veins
        \item Assessment of venous drainage
        \item Optimal for parenchymal volume measurements
    \end{itemize}
\end{itemize}

\subsubsection{Late Phase (Delayed Phase)}
\begin{itemize}
    \item \textbf{Timing}: 3-5 minutes after contrast injection
    \item \textbf{Purpose}: Captures contrast excretion and collecting system
    \item \textbf{Key Features}:
    \begin{itemize}
        \item Contrast in renal collecting system
        \item Assessment of kidney function
        \item Visualization of ureters and bladder
        \item Evaluation of contrast excretion
    \end{itemize}
\end{itemize}

\subsection{Technical Specifications}
\begin{itemize}
    \item \textbf{Image Format}: DICOM (Digital Imaging and Communications in Medicine)
    \item \textbf{File Structure}: Organized by case numbers (1-25)
    \item \textbf{Data Organization}:
    \begin{itemize}
        \item \texttt{IOD\_arterial/}: Arterial phase DICOM files
        \item \texttt{IOD\_venous/}: Venous phase DICOM files  
        \item \texttt{IOD\_late/}: Late phase DICOM files
    \end{itemize}
\end{itemize}

\section{Kidney Segmentation and Analysis}

\subsection{Segmentation Process}

The kidney segmentation process involves precise delineation of kidney structures from CT images using advanced image processing techniques.

\subsubsection{Segmented Structures}
\begin{enumerate}
    \item \textbf{Kidney Parenchyma}: The functional tissue of the kidney
    \item \textbf{Renal Arteries}: Blood vessels supplying the kidney
    \item \textbf{Renal Veins}: Blood vessels draining the kidney
    \item \textbf{Cortex}: Outer layer of kidney tissue
    \item \textbf{Medulla}: Inner region of kidney tissue
\end{enumerate}

\subsubsection{Segmentation Output Files}
\begin{itemize}
    \item \textbf{\texttt{.seg.nrrd}}: Segmentation files in NRRD format
    \item \textbf{\texttt{.nrrd}}: Volume data files
    \item \textbf{\texttt{.csv}}: Measurement tables with quantitative data
    \item \textbf{\texttt{.mrml}}: 3D Slicer scene files for visualization
\end{itemize}

\subsubsection{Quantitative Measurements}
The segmentation provides detailed measurements including:
\begin{itemize}
    \item \textbf{Volume measurements} (mm\textsuperscript{3}, cm\textsuperscript{3})
    \item \textbf{Voxel counts} for each segmented structure
    \item \textbf{Intensity statistics} (mean, median, standard deviation)
    \item \textbf{Percentile values} (5th, 95th percentiles)
    \item \textbf{Minimum and maximum intensity values}
\end{itemize}

\subsection{Example Measurement Data}
\begin{verbatim}
Segment: arterial_threshold_1
- Volume: 1,639.82 cm³
- Mean intensity: 194.116 HU
- Standard deviation: 98.376 HU
- Voxel count: 2,944,174
\end{verbatim}

\section{eGFR (Estimated Glomerular Filtration Rate) Calculations}

\subsection{What is eGFR?}

eGFR is a crucial measure of kidney function that estimates how well the kidneys are filtering waste from the blood. It's calculated using serum creatinine levels and other factors.

\subsection{Clinical Significance}
\begin{itemize}
    \item \textbf{Normal eGFR}: $>$90 mL/min/1.73m\textsuperscript{2}
    \item \textbf{Mild reduction}: 60-89 mL/min/1.73m\textsuperscript{2}
    \item \textbf{Moderate reduction}: 30-59 mL/min/1.73m\textsuperscript{2}
    \item \textbf{Severe reduction}: 15-29 mL/min/1.73m\textsuperscript{2}
    \item \textbf{Kidney failure}: $<$15 mL/min/1.73m\textsuperscript{2}
\end{itemize}

\subsection{Data Structure in This Project}

The eGFR data is stored in CSV files with the following structure:
\begin{itemize}
    \item \textbf{Patient ID}: Unique identifier for each case
    \item \textbf{Date}: Date of eGFR measurement
    \item \textbf{eGFR Value}: Calculated eGFR in mL/min/1.73m\textsuperscript{2}
    \item \textbf{Serum Creatinine}: Laboratory value used in calculation
    \item \textbf{Measurement Type}: Method used for eGFR calculation
\end{itemize}

\subsection{Longitudinal Analysis}
The project includes longitudinal eGFR data showing:
\begin{itemize}
    \item \textbf{Multiple measurements} over time for each patient
    \item \textbf{Trend analysis} of kidney function
    \item \textbf{Correlation} between imaging findings and functional decline
\end{itemize}

\section{Project Data Structure}

\subsection{Case Organization}
\begin{verbatim}
Cases/
├── 1/                    # Case number
│   ├── DICOM/           # Original DICOM files
│   │   ├── IOD_arterial/ # Arterial phase images
│   │   ├── IOD_venous/   # Venous phase images
│   │   └── IOD_late/     # Late phase images
│   ├── eGFR_1.csv       # eGFR measurements
│   └── Segmenteringer/  # Segmentation results
│       ├── *.seg.nrrd   # Segmentation files
│       ├── *.nrrd       # Volume files
│       ├── *.csv        # Measurement tables
│       └── *.mrml       # 3D Slicer scenes
\end{verbatim}

\subsection{Key File Types}
\begin{itemize}
    \item \textbf{DICOM}: Medical imaging standard format
    \item \textbf{NRRD}: Nearly Raw Raster Data format for volumes
    \item \textbf{SEG.NRRD}: Segmentation data format
    \item \textbf{CSV}: Tabular data for measurements
    \item \textbf{MRML}: 3D Slicer scene files
\end{itemize}

\section{Clinical Applications}

\subsection{Diagnostic Capabilities}
\begin{enumerate}
    \item \textbf{Kidney Function Assessment}: Quantitative evaluation of kidney performance
    \item \textbf{Disease Progression}: Longitudinal monitoring of kidney function
    \item \textbf{Treatment Planning}: Data-driven decisions for patient care
    \item \textbf{Research Applications}: Large-scale analysis of kidney health
\end{enumerate}

\subsection{Research Value}
\begin{itemize}
    \item \textbf{Population Studies}: Analysis across multiple patients
    \item \textbf{Method Validation}: Comparison of different measurement techniques
    \item \textbf{Clinical Outcomes}: Correlation between imaging and functional data
    \item \textbf{Biomarker Development}: Identification of imaging biomarkers
\end{itemize}

\section{Technical Tools and Software}

\subsection{Primary Software}
\begin{itemize}
    \item \textbf{3D Slicer}: Open-source medical imaging platform
    \item \textbf{Python}: Data analysis and processing
    \item \textbf{Jupyter Notebooks}: Interactive analysis environment
\end{itemize}

\subsection{Data Analysis Capabilities}
\begin{itemize}
    \item \textbf{Volume Calculations}: Precise kidney volume measurements
    \item \textbf{Intensity Analysis}: Quantitative assessment of contrast enhancement
    \item \textbf{Statistical Analysis}: Comprehensive data processing
    \item \textbf{Visualization}: 3D rendering and multi-planar views
\end{itemize}

\section{Future Directions}

\subsection{Potential Enhancements}
\begin{enumerate}
    \item \textbf{Machine Learning}: Automated segmentation and analysis
    \item \textbf{Predictive Modeling}: Risk stratification based on imaging data
    \item \textbf{Standardization}: Development of standardized protocols
    \item \textbf{Integration}: Connection with electronic health records
\end{enumerate}

\subsection{Research Opportunities}
\begin{itemize}
    \item \textbf{Biomarker Discovery}: Identification of new imaging biomarkers
    \item \textbf{Population Health}: Large-scale epidemiological studies
    \item \textbf{Precision Medicine}: Personalized treatment approaches
    \item \textbf{Clinical Trials}: Support for interventional studies
\end{itemize}

\bigskip
\hrule
\bigskip

\textit{This project represents a comprehensive approach to kidney analysis using modern medical imaging techniques, combining structural assessment through CT segmentation with functional evaluation through eGFR calculations.}

\end{document}
